% A self-contained section to insert into the Gomez--Venegas draft.
% IMPORTANT: G result relies on the center theorem for h=1 in the supplied manuscript.
\section{The ozone group of $\mathbb G(p,f,1)$}
\label{sec:Gozcomplete}
\begin{theorem}\label{thm:Gozcomplete}
Let $\Bbbk$ be algebraically closed of characteristic zero, let $p$ be a
primitive $n$-th root of unity with $n\ge3$, and let $f\in\Bbbk^\times$
be an admissible parameter of the family $\mathbb G$. Put
$k=n/2$ when $n$ is even. Then
\[
 \operatorname{Oz}(\mathbb G(p,f,1))\cong
 \begin{cases}
 (\mathbb Z/n)^2,&n\text{ odd or }n\equiv2\pmod4,\\
 \mathbb Z/(n/2)\times\mathbb Z/2\times\mathbb Z/n,&4\mid n.
 \end{cases}
\]
In particular $|\operatorname{Oz}(\mathbb G(p,f,1))|=n^2$.
\end{theorem}
\begin{proof}
Write $X=x_2+x_1/(p-1)$. The defining relations in this basis are
\begin{align*}
 x_1X&=Xx_1,&y_2y_1&=py_1y_2,\\
 y_1x_1&=px_1y_1,&y_2x_1&=px_1y_2,\\
 y_1X&=p^2Xy_1+x_1y_2,&y_2X&=f x_1y_1+Xy_2.
\end{align*}
Every ozone automorphism is bigraded (Lemma~\ref{lem:ozbig}). The center
calculation in Theorem~\ref{thm:G} shows that $x_1^n,y_1^n,y_2^n$ belong
to the center. As $x_1,X$ commute, fixing $x_1^n$ implies
$\sigma(x_1)=a x_1$ with $a^n=1$. In the quantum $p$-plane, the
root-of-unity binomial identity
$(r y_1+s y_2)^n=r^n y_1^n+s^n y_2^n$
implies that fixing $y_1^n,y_2^n$ separately forces
$\sigma(y_1)=c y_1$ and $\sigma(y_2)=d y_2$, with $c^n=d^n=1$.
Write $\sigma(X)=b x_1+t X$. The relation $y_1X=p^2Xy_1+x_1y_2$
gives $bc(p-p^2)=0$, so $b=0$, and $tc=ad$. The relation
$y_2X=f x_1y_1+Xy_2$ gives $td=ac$. Thus
\[
 t=\varepsilon a,\qquad d=\varepsilon c,\qquad
 \varepsilon\in\{1,-1\}.
\]
Conversely these maps preserve all six defining relations.

For odd $n$, Theorem~\ref{thm:G}(a) says the center is generated by
$x_1^n,y_1^n,y_2^n$ and
$u_n=D_n(X,\alpha x_1^2)$, where
$\alpha=-f/(p^2-1)^2$ and $D_n$ is the Dickson polynomial.
Every monomial in $u_n$ has the form $X^{n-2j}x_1^{2j}$;
therefore $\sigma(u_n)=\varepsilon^n a^n u_n$. It is fixed if and
only if $\varepsilon=1$. We obtain $\mu_n^2$.

For even $n=2k$, Theorem~\ref{thm:G}(b) gives the center generators
$x_1^n,y_1^n,y_2^n,u_k=D_k(X,\alpha x_1^2)$ and
$w=x_1^k y_1^k y_2^k$. Their multipliers under $\sigma$ are,
respectively, $a^n,c^n,(\varepsilon c)^n,
(\varepsilon a)^k$, and $a^k(c\varepsilon c)^k$.
Thus all generators are fixed exactly when
\[
 a^n=c^n=1,\qquad (\varepsilon a)^k=1.
\]
The last equation also fixes $w$ since its multiplier is
$a^k\varepsilon^kc^{2k}=1$.
If $k$ is odd, $\varepsilon=a^{-k}$ is uniquely determined by
$a\in\mu_n$, so the group is $\mu_n^2$. If $k$ is even, the equation
forces $a\in\mu_k$ and leaves $\varepsilon\in\mu_2$ free; hence the
group is $\mu_k\times\mu_2\times\mu_n$. The maps form finite groups
because their action on the four generators is diagonal.
\end{proof}

\begin{remark}
The theorem is for $h=1$ only. For general roots of unity $h$,
Theorem~\ref{thm:G} does not describe the full twisted center, and the
conclusion must NOT be extrapolated. In particular the table already
shows a change of the ozone group at $(p,h)=(i,i)$.
\end{remark}

\subsection{What the same argument gives for $\mathbb D$}
\begin{proposition}\label{prop:Dnecessary}
Let $p$ be a primitive $n$-th root of unity with $n\ge3$ and put
$e=\mathrm{lcm}(n,2)$. For $B=\mathbb D(p,1)$ every ozone automorphism
has the form
\[
 x_1\mapsto a x_1,\quad x_2\mapsto\varepsilon a x_2,\quad
 y_1\mapsto c y_1,\quad y_2\mapsto\varepsilon c y_2,
\]
where $a^e=c^e=1$, $\varepsilon^2=1$, and
$(\varepsilon c)^n=1$. These are necessary conditions, not an
exhaustive description of $\operatorname{Oz}(B)$.
\end{proposition}
\begin{proof}
The central elements $x_1^e,y_1^e,y_2^n$ are constructed in the
proof of Theorem~\ref{thm:viaA}. Bigradedness and the identity
$(r x_1+s x_2)^2=r^2x_1^2+s^2x_2^2$ in the $(-1)$-quantum plane
show that fixing $x_1^e$ forces $\sigma(x_1)=a x_1$ with $a^e=1$.
Fixing $y_1^e$ and $y_2^n$ in the quantum $p$-plane similarly
forces $\sigma(y_1)=c y_1$ and $\sigma(y_2)=d y_2$ with
$c^e=d^n=1$. The $(-1)$-plane relation makes
$\sigma(x_2)=t x_2$. Both nontrivial mixed relations give
$tc=ad$ and $td=ac$. These imply $t=\varepsilon a$ and
$d=\varepsilon c$, as stated.
\end{proof}
\begin{remark}
For $n=4$, the necessary conditions admit 32 maps, whereas the
existing example table reports $|\operatorname{Oz}(\mathbb D(i,1))|=16$.
Thus the central elements used here do NOT suffice for exhaustivity:
additional generators of the full center (and their action) are
required. This is a concrete, testable obstruction to closing $\mathbb D$.
\end{remark}
